% s-343-p343-m002; original PDF 343, 344
\Aphorism{s-343-p343-m002}{395}
Noch einige Stufen aufwärts im Erkennen der Seele, und es tagt uns, dass unser zuversichtliches „Ich kann thun, was ich will“ dieselbe Einfalt und Komödie ist, wie wenn man frühmorgens der Sonne aufzugehen geböte. Der „Wille“ ist eine Begleiterscheinung, eine beiläufige Auskunft, die uns über uns selbst ertheilt, nicht ein Gutachten, das uns abverlangt wird; die willkürliche Handlung unterscheidet sich von der unwillkürlichen nur durch den Geschäftsgang, dadurch dass das Bewusstsein eine Abschrift des geheimen Befehls empfängt — aber das Bewusstsein selbst hat nichts zu befehlen. Das Bewusstsein ist die Redaction, die von officiöser Seite auf dem Laufenden erhalten wird, aber nichts zu beantragen, zu beschliessen oder zu verantworten hat. Wir erfahren von dem, was in uns und ausser uns vorgeht, nur in der Form des Bewusstseins, durch die Berichterstattung des Bewusstseins —: das verführt uns zu glauben, das Bewusstsein ziehe auch die Fäden und drehe die Räder. Die Schopenhauersche Allmacht des Willens ist, was immer sonst, ein Abweg vom heutigen Gange des Erkennens; wir werden vielleicht noch in diesem Jahrhundert die Ohnmacht des Willens beweisen und ihn, die missverstandenste aller metaphysischen Erschleichungen, in sein Nichts zurückverweisen.
\par

